% ---
% id: protocol-compiler/rem:alphabet
% kind: rem
% title: "The outcome alphabet is chosen upstream"
% summary: >-
%   Bob's emitted outcome alphabet is chosen upstream, the receiver's declared
%   coarse-graining deciding how much of the click pattern survives. Every
%   object of this module indexed by it is stated over any finite alphabet, so a
%   fine-grained instance changes declared data and no claim here.
% prominence: standard
% status: accepted
% uses: [def:declaration, def:spec, def:chaincond, def:phi, def:weighted]
% mentions: [rem:coarse, rem:relabel, receiver/def:clickalphabet, receiver/def:coarse, receiver/thm:coarse-graining, reductions/def:joint-family]
% assumes: []
% refs: []
% provenance: newly-derived
% lean:
% ---

\begin{remark}\label{rem:alphabet}
Nothing in this module fixes what Bob's \tref{outcome alphabet} is. The emitted
$\Vset$ is the alphabet of the family the last link of $\Chrec$ delivered
(Definition~\ref{def:spec}\,(c)), and that family's alphabet is the upstream
modules': the raw outcome of a threshold-detector receiver is the
\emph{click pattern}, the set $S\subseteq[K]$ of ports that fired
(\xref{receiver/def:clickalphabet}), and it is the receiver's
\tref{declared coarse-graining} --- a surjection from $2^{[K]}$ onto the
emitted alphabet (\xref{receiver/def:coarse},
\xref{receiver/thm:coarse-graining}) --- that decides how much of the pattern
survives, \emph{before} the joint alphabet with Bob's setting is formed
(\xref{reductions/def:joint-family}) and before any reduction runs. A
coarse-graining that keeps every pattern, one that keeps only click, no-click
and multi-click, and one that keeps the single clicks and merges the rest are
three declarations of that map, and this module receives the result of
whichever was declared.

Every object of this module that is indexed by $\Vset$ is stated over an
arbitrary finite non-empty alphabet and admits any finite disjoint union.
The family slot $\{F_v\}_{v\in\Vset}$ of \eqref{eq:declaration} is a family
of block lists over a finite alphabet, and D(iii) is decided element by
element whatever $|\Vset|$ is. The outcome correspondence $\gamma$ of
Definition~\ref{def:chaincond}\,(C2) is a total surjection from any finite
physical alphabet onto any finite emitted one, the direction being the only
thing (C2) fixes (Remark~\ref{rem:coarse}), so a finer emitted alphabet is a
$\gamma$ with smaller fibres and nothing else. The readouts of the
\tref{party partition} and Bob's \tref{announcement map} $a_B$
(Definition~\ref{def:spec}\,(d))
are declared on the alphabets the specification names, and $\phi$
(Definition~\ref{def:phi}) on the announcement classes; a finer $\Vset$ is a
larger table for each and no new condition. And both forms of
Definition~\ref{def:weighted} are constructions over a disjoint union ---
$\bigsqcup_t\Uset_t\times\{t\}$, and the support $\Vset_P\subseteq
\Vset\times\mathcal{F}$ --- so a disjoint union over settings and choice
values is what they produce, not something they have to be extended to.

A \emph{fine-grained instance} --- one whose receiver keeps more of the click
pattern than the coarsest sound choice, or one whose outcome carries a
setting and a choice value beside the detection result --- therefore changes
the receiver's declared coarse-graining and the instance's public part, and
changes no claim of this module: the compiler's claims are stated for the
alphabet it is handed. What it does change downstream is not this module's
either: the size of the declaration, the number of exact decisions V3 takes,
and which coarse outcomes carry a bound
(\xref{receiver/thm:coarse-graining}\,(iii): a bound on a coarse outcome does
not descend to its fibres), which is the receiver's trade-off and the reason
its coarse-graining is declared data rather than a convention. Which
alphabets that coarse-graining can reach is the receiver's question; where a
setting-dependent choice value enters the alphabet is
Remark~\ref{rem:relabel}.
\end{remark}
